% ====================================================================
% VİZE RAPOR ŞABLONU – Türk Masalları Derlemesi Web Sayfası
% --------------------------------------------------------------------
% DERLEME: Overleaf'te Menu > Compiler > XeLaTeX (pdfLaTeX de çalışır).
% ====================================================================
\documentclass[12pt,a4paper]{article}

\usepackage{iftex}
\ifPDFTeX
  \usepackage[T1]{fontenc}
  \usepackage[utf8]{inputenc}
\else
  \usepackage{fontspec}
\fi
\usepackage[turkish]{babel}
\usepackage[margin=2.4cm]{geometry}
\usepackage{setspace}
\usepackage{graphicx}
\usepackage{booktabs}
\usepackage{longtable}
\usepackage{array}
\usepackage{colortbl}
\usepackage{xcolor}
\usepackage{enumitem}
\usepackage{listings}
\usepackage{amssymb}
\usepackage{caption}
\usepackage{fancyhdr}
\usepackage{hyperref}

\onehalfspacing
\setlength{\parindent}{0pt}
\setlength{\parskip}{6pt}
\setlist{itemsep=2pt, topsep=4pt}
\hypersetup{colorlinks=true, linkcolor=blue!50!black, urlcolor=blue!50!black}

\definecolor{ana}{HTML}{14325C}
\definecolor{acik}{HTML}{EEF2F8}

\lstset{basicstyle=\ttfamily\small, backgroundcolor=\color{acik},
  frame=single, rulecolor=\color{ana}, breaklines=true,
  columns=fullflexible, keepspaces=true, showstringspaces=false}

\pagestyle{fancy}
\setlength{\headheight}{14.5pt} 
\fancyhf{}
\fancyhead[L]{Türk Masalları Derlemesi}
\fancyhead[R]{Vize Raporu}
\fancyfoot[C]{\thepage}

% ---------- Yardımcı komutlar ----------
\newcommand{\haftagorsel}[1]{%
  \IfFileExists{#1}{%
    \includegraphics[width=\linewidth,height=4.2cm,keepaspectratio]{#1}%
  }{%
    \fbox{\parbox[c][2.6cm][c]{\dimexpr\linewidth-2\fboxsep-2\fboxrule\relax}{%
      \centering\small\color{gray} Ekran görüntüsü / görsel ekleyin\\ \texttt{#1}}}%
  }%
}
\newcommand{\haftaek}[1]{%
  \IfFileExists{#1}{\\[4pt]\includegraphics[width=\linewidth,height=4.2cm,keepaspectratio]{#1}}{}%
}

% ---------- Kapak bilgileri ----------
\newcommand{\universite}{İstanbul Medipol Üniversitesi / Görsel İletişim Tasarımı}
\newcommand{\dersadi}{Programlamaya Giriş}
\newcommand{\hocaadi}{Meral Şentürk}
\newcommand{\ogrenciadi}{Ozan Ali Bugaoğlu}
\newcommand{\ogrencino}{86230065}
\newcommand{\projeadi}{Masal Kutusu}
% Bağlantıları tam adres olarak yazın (https:// ile başlayarak)
\newcommand{\githublink}{\url{https://github.com/ozanbuga/masal-kutusu}}
\newcommand{\youtubelink}{\url{https://www.youtube.com/watch?v=sknqtlP4rMM}}

\begin{document}

% ====================================================================
% KAPAK
% ====================================================================
\begin{titlepage}
\centering
\vspace*{1cm}
{\large \universite \par}
\vspace{1.8cm}
{\Huge\bfseries\color{ana} \projeadi \par}
\vspace{0.4cm}
{\Large Türk Masalları Derlemesi Web Sayfası \par}
{\large Vize Raporu \par}
\vspace{1.8cm}
{\large \textbf{Ders:} \dersadi \par}
{\large \textbf{Öğretim Üyesi:} \hocaadi \par}
\vspace{1.2cm}
{\large \textbf{Hazırlayan:} \ogrenciadi \par}
{\large \textbf{Öğrenci No:} \ogrencino \par}
\vspace{1.8cm}
\fbox{\parbox{0.88\linewidth}{\centering\vspace{6pt}
{\large\bfseries Proje Bağlantıları}\\[8pt]
\textbf{GitHub deposu:}\\ \githublink \\[8pt]
\textbf{YouTube sunum videosu:}\\ \youtubelink
\vspace{6pt}}}
\vfill
{\large \today \par}
\end{titlepage}

\tableofcontents
\newpage

% ====================================================================
\section{Giriş}
% ====================================================================
Bu projenin amacı, Türk halk anlatı geleneğinin önemli bir parçası olan masalları, modern web tasarım (UI/UX) standartlarıyla dijital bir okuma deneyimine dönüştürmektir. Proje, yönergedeki katı teknik kısıtlamalara sadık kalarak, hiçbir harici kütüphane veya çerçeve (framework) kullanılmadan geliştirilmiştir. İstenilen 8 masal, tek bir HTML dosyası (\texttt{tek.html}) içerisinde toplanmış; sitenin tüm stil yönergeleri \texttt{<style>} etiketleriyle, interaktif algoritmaları ise \texttt{<script>} etiketleriyle aynı dosya içine gömülmüştür. Görseller harici sunuculardan çekilmek yerine yerel dizinde tutularak, projenin tamamen çevrimdışı (offline) çalışabilmesi sağlanmıştır.

\subsection{Proje Özeti}
"Masal Kutusu", okuyuculara 8 farklı anonim masalı anında okuma, arama ve filtreleme imkanı sunan; gece modu ve dinamik metin boyutu ayarlama özelliklerine sahip, mobil uyumlu (responsive) bir tek sayfa web sitesidir (Single Page Application).

\subsection{Gereksinimlerin Karşılanması}
\begin{center}
\rowcolors{2}{acik}{white}
\begin{tabular}{@{}p{8.2cm}c@{}}
\toprule
\rowcolor{ana}\textcolor{white}{\textbf{Gereksinim}} & \textcolor{white}{\textbf{Durum}} \\
\midrule
Masal sayısı 8--12 arasında (8 masal)         & $\boxtimes$ \\
Tek HTML dosyası                                    & $\boxtimes$ \\
Stil kodları içsel (\texttt{<style>} içinde)        & $\boxtimes$ \\
Betik kodları içsel (\texttt{<script>} içinde)      & $\boxtimes$ \\
İnternetsiz açılıp çalışıyor                        & $\boxtimes$ \\
JavaScript sayfada anlamlı bir işlev görüyor        & $\boxtimes$ \\
Masal metinleri kendi cümlelerimle yazıldı          & $\boxtimes$ \\
Sayfa telefonda da düzgün görünüyor                 & $\boxtimes$ \\
\bottomrule
\end{tabular}
\end{center}

% ====================================================================
\section{Masalların Seçimi}
% ====================================================================
\begin{longtable}{@{}cp{4cm}p{3cm}p{5.4cm}@{}}
\toprule
\textbf{No} & \textbf{Masal} & \textbf{Tür} & \textbf{Seçilme nedeni} \\
\midrule
\endhead
1 & Keloğlan ve Sihirli Tas & Anadolu Masalı & Türk halk kültürünün en bilinen zeka ve şans motifini barındırması. \\
2 & Billur Köşk & Saray Masalı & Fantastik ögeler ve bilmece çözme mantığıyla sürükleyici bir anlatı sunması. \\
3 & Nardaniye Hanım & Anadolu Masalı & Pamuk Prenses masalının çok daha eski bir Anadolu varyasyonu olması. \\
4 & Tembel Ahmet & Güldürücü Masal & Mizahi bir dille çalışkanlığın önemini vurgulayan eğlenceli bir yapıya sahip olması. \\
5 & Yedi Başlı Ejderha & Kahramanlık Mas. & Mitolojik ögelerle destansı bir okuma deneyimi sunması. \\
6 & Sabır Taşı & Anadolu Masalı & Derin bir duygusal temaya ve kültürel sabır konseptine odaklanması. \\
7 & Üç Turunçlar & Olağanüstü Masal & Büyü ve sihir ögeleriyle masal diyarının fantastik sınırlarını yansıtması. \\
8 & Altın Bülbül & Hayvan Masalı & İnsan ve hayvan yardımlaşmasını ön plana çıkaran geleneksel bir kurgu olması. \\
\bottomrule
\end{longtable}

% ====================================================================
% ====================================================================
\section{Tasarım Kararları}
% ====================================================================
Bu bölümde, projenin görsel ve işlevsel altyapısını oluştururken alınan temel tasarım kararları, renk paleti seçimleri, tipografi tercihleri ve kullanıcı deneyimi (UX) yaklaşımları detaylandırılmaktadır.

% ====================================================================

\subsection{Konuya Uygunluk}
Arayüz tasarımı, modern ve ruhsuz bir web sitesi yerine eski basılı bir masal kitabını andıracak şekilde kurgulanmıştır. Sayfa zemininde saf beyaz yerine gözü yormayan eski kağıt tonları kullanılmış; başlık altlarına saf CSS ile 45 derece döndürülmüş kırmızı kareler çizilerek geleneksel Anadolu kilim motiflerine atıfta bulunulmuştur. Paragrafların ilk harfleri "Drop Cap" tekniğiyle büyütülerek tematik bütünlük sağlanmıştır.

\subsection{Renk Paleti ve Yazı Tipi}
\begin{center}
\rowcolors{2}{acik}{white}
\begin{tabular}{@{}llp{6.4cm}@{}}
\toprule
\rowcolor{ana}\textcolor{white}{\textbf{Ad}} & \textcolor{white}{\textbf{Kod}} & \textcolor{white}{\textbf{Nerede / neden kullandınız?}} \\
\midrule
Kağıt Kremi & \texttt{\#f9f7f1} & Arka planda, eski kitap sayfası hissi vermek için. \\
Kök Kırmızısı & \texttt{\#8a2a2a} & Başlıklarda, butonlarda ve ayraçlarda (Anadolu vurgusu). \\
Koyu Füme & \texttt{\#2b2b2b} & Ana metinlerde, sert kontrasttan kaçınmak için. \\
Serif Yazı & \texttt{Georgia} & Ana okuma alanında, satır takibini kolaylaştırdığı için. \\
\bottomrule
\end{tabular}
\end{center}

\subsection{Okunabilirlik, Gezinme ve Duyarlı Tasarım}
\begin{itemize}
  \item \textbf{Okunabilirlik:} Metin satır uzunluğu maksimum 720px ile sınırlandırılarak gözün kaybolması engellenmiştir. Satır aralığı (line-height) 1.8'dir.
  \item \textbf{Gezinme:} Ekranın soluna sabitlenen (sticky) ve arama yapılabilen bir menü yerleştirilmiş, farklı sayfalara gitmek yerine JavaScript ile aynı sayfada sekme geçişi sağlanmıştır.
  \item \textbf{Duyarlı tasarım:} \texttt{Flexbox} mimarisiyle ekran 768px altına düştüğünde menü otomatik olarak sayfanın en üstüne yerleşmektedir.
  \item \textbf{Tutarlılık:} Gece modu dahil tüm renk ve ölçüler \texttt{:root} içerisindeki CSS değişkenlerinden çekilmiştir.
\end{itemize}

\subsection{Ekran Görüntüleri}
\shorthandoff{=}
\begin{figure}[h!]
  \centering
  \includegraphics[width=0.9\textwidth]{ekran-masaustu.png}
  \caption{Sayfanın bilgisayar görünümü}
\end{figure}
\begin{figure}[h!]
  \centering
  \includegraphics[width=0.35\textwidth]{ekran-mobil.png}
  \caption{Sayfanın telefon görünümü}
\end{figure}
\clearpage
% ====================================================================
\section{Gerçekleştirme}
% ====================================================================
Bu proje, kullanıcı deneyimini kesintiye uğratmamak adına Çoklu Sayfa (Multi-Page) mimarisi yerine Tek Sayfa Uygulaması (Single Page Application - SPA) mantığıyla geliştirilmiştir. Kullanıcı sol menüden bir masal seçtiğinde sayfa yenilenmez; JavaScript'in DOM manipülasyonu yetenekleri kullanılarak yalnızca ilgili masalın görünürlüğü CSS sınıfı eklenerek (\texttt{active}) açılırken, diğer masallar gizlenir. HTML semantiği, CSS Flexbox yapısı ve JavaScript olay dinleyicileri (event listeners) tek bir dosyada kusursuz bir uyum içinde çalışarak hızlı ve tamamen çevrimdışı (offline) bir okuma deneyimi sunmaktadır.

\subsection{HTML Yapısı}
HTML kodlaması, div çöplüğü yaratmadan HTML5'in semantik etiketleriyle yapılmıştır. Tepe kontrolleri \texttt{<header>}, yan menü \texttt{<aside>}, masal okuma alanı \texttt{<main>} ve her bir masal \texttt{<article>} etiketiyle sarmalanmıştır. Liste elemanlarına JS için \texttt{data-isim} ve \texttt{data-tur} nitelikleri atanmıştır.

\subsection{İçsel CSS}
Arayüz düzeni Flexbox yapısıyla kurgulanmış, karanlık mod (dark mode) için CSS değişkenleri kullanılmıştır.

\begin{lstlisting}[caption={Icsel CSS'ten onemli bir bolum}]
.ayrac::after {
    content: "";
    width: 12px;
    height: 12px;
    background-color: var(--primary-color);
    transform: rotate(45deg); /* Kareyi eskenar dortgene cevirme */
    z-index: 1;
}
.masal-icerik p:first-of-type::first-letter {
    float: left;
    font-size: 4rem;
    font-weight: bold;
    color: var(--primary-color);
}
\end{lstlisting}

\subsection{İçsel JavaScript}
Masallar farklı sayfalara bölünmek yerine, DOM manipülasyonu ile yönetilmiştir. \texttt{masallariFiltrele} fonksiyonu, kullanıcının girdiği harfleri dinleyerek anlık arama yapar.

\begin{lstlisting}[caption={Icsel JavaScript'ten onemli bir bolum}]
function masallariFiltrele() {
    const aranan = document.getElementById('aramaKutusu').value.toLowerCase();
    const tur = document.getElementById('kategoriFiltresi').value;
    
    document.querySelectorAll('#masalListesi li').forEach(madde => {
        const isimUyuyor = madde.getAttribute('data-isim').includes(aranan);
        const turUyuyor = (tur === "Tumu" || tur === madde.getAttribute('data-tur'));
        
        if (isimUyuyor && turUyuyor) madde.classList.remove('gizli');
        else madde.classList.add('gizli');
    });
}
\end{lstlisting}

\subsection{Sayfanın Özellikleri}
JavaScript sayfada yalnızca süs amacıyla değil, anlamlı işlevler (UX) sunmak için kullanılmıştır. Sayfada yer alan temel dinamik özellikler şunlardır:
\begin{itemize}
    \item \textbf{Gece Modu (Dark Mode):} \texttt{geceModuGecis()} fonksiyonu ile \texttt{body} etiketine \texttt{dark-mode} sınıfı eklenir/çıkarılır. Böylece tüm renk paleti anında değişir.
    \item \textbf{Dinamik Metin Boyutu:} \texttt{degistirMetinBoyutu()} fonksiyonu, CSS \texttt{:root} içerisindeki \texttt{--font-size-base} değişkenini JS ile manipüle ederek, sayfayı yenilemeden tüm yazıları eşzamanlı büyütür veya küçültür.
    \item \textbf{Okuma İlerleme Çubuğu:} Sayfa kaydırıldıkça (scroll) çalışan olay dinleyicisi (event listener), matematiksel olarak dikey ilerleme oranını hesaplayarak sayfa üstündeki kırmızı çubuğun genişliğini (\texttt{width}) yüzde olarak doldurur.
\end{itemize}

% ====================================================================
\section{Yapay Zekâ Kullanımı}
% ====================================================================

\begin{longtable}{@{}p{2.6cm}p{3.6cm}p{4.2cm}p{3.4cm}@{}}
\toprule
\textbf{Araç} & \textbf{Ne için kullandınız?} & \textbf{İstem (prompt) özeti} & \textbf{Sayfada nerede?} \\
\midrule
\endhead
Google Gemini & Metin Derleme (Metin) & "Anonim Türk masallarını 2 dakikalık okuma süresine göre akıcı dille özetle." & Okuma alanı (article) \\
Google Gemini & Görsel Üretimi (Görsel) & İsmini vereceğim masallara göre illustrastif görseller üret. & Tüm masal içeriklerindeki \texttt{<img>} etiketlerinde \\
Google Gemini & Kod Analizi (Kod) & "Görselleri Base64'e çevirmek vs klasörde tutmak arasındaki mimari farklar nelerdir?" & Proje dosya mimarisi \\
Google Gemini & Kod Desteği (Kod) & "Arama ve kategori filtresini aynı fonksiyonda mantıksal operatörle nasıl yazarım?" & Script bloğu \\
\bottomrule
\end{longtable}

\subsection{Anladığım ve Kendim Yaptığım Kısımlar}
Yapay zeka asistanı yönergedeki "tek dosya" kuralını ihlal etmemek adına görselleri Base64 formatına çevirmeyi önermiş, ancak bu durumun Clean Code yapısını bozacağını analiz ederek inisiyatif aldım. Resimleri yönergedeki internetsiz çalışma kuralını mükemmel sağlayan yerel dizin mantığıyla kendim bağladım. Flexbox medya sorgularını ve DOM elementlerine \texttt{active} / \texttt{gizli} sınıfı ekleme algoritmalarını tasarım disiplinime uygun şekilde kendi kurgumla revize ettim.

% ====================================================================
\section{Haftalık Katkılar}
% ====================================================================

\begin{longtable}{@{}>{\raggedright\arraybackslash}p{1.7cm}%
                    >{\raggedright\arraybackslash}p{6.4cm}%
                    >{\centering\arraybackslash}p{6cm}@{}}
\toprule
\rowcolor{ana}\textcolor{white}{\textbf{Hafta}} & \textcolor{white}{\textbf{Yapılan çalışma}} & \textcolor{white}{\textbf{Görsel}} \\
\midrule
\endfirsthead
\toprule
\rowcolor{ana}\textcolor{white}{\textbf{Hafta}} & \textcolor{white}{\textbf{Yapılan çalışma}} & \textcolor{white}{\textbf{Görsel}} \\
\midrule
\endhead
\bottomrule
\endfoot

\textbf{1. hafta} &
Bu haftadaki görevleri sonradan telafi ettim. hazır HTML taslağında CSS kullanarak arka plan rengini maviye, yazı rengini beyaza ve fontu Verdana olarak değiştirdim. HTML yapısını koruyarak CSS ile görsel düzenleme yapmayı öğrendim.\newline \textbf{Commit:} CSS ile mavi tema ve tipografi güncellemeleri uygulandı
\par\vspace{3pt}\small $\square$ Ders içinde katılım sergiledim\par $\boxtimes$ Telafi yaptım &
\haftagorsel{hafta1.png}\haftaek{hafta1b.png} \\
\midrule
\textbf{2. hafta} &
Bu haftadaki görevleri sonradan telafi ettim. Temel HTML etiketleri (h1, h2, p, ul, a) ve içsel CSS kullanarak "Evcil Kediler" temalı tek sayfalık bir rehber hazırladım. Sayfa içi bağlantı (anchor link) özelliği ekleyerek, menüdeki başlıklara tıklandığında ilgili içerik bölümüne gidilmesini sağladım. Renk ve metin stillerini CSS ile özelleştirdim.\newline \textbf{Commit:} Evcil Kediler temel HTML ve CSS tasarımı eklendi
\par\vspace{3pt}\small $\square$ Ders içinde katılım sergiledim\par $\boxtimes$ Telafi yaptım &
\haftagorsel{hafta2.png}\haftaek{hafta2b.png} \\
\midrule
\textbf{3. hafta} &
Bu haftadaki görevleri sonradan telafi ettim. HTML, CSS (Grid ve Flexbox) ile üç sayfalık (Platform, Kayıt Ol, Giriş Yap) bir yayın platformu arayüzü tasarladım. Form elemanları ekledim ve JavaScript kullanarak form doğrulama (boş alan kontrolü) ile şifre alanını "Göster/Gizle" işlevlerini kodladım.\newline \textbf{Commit:} Yayın platformu arayüzü, formlar ve JS şifre göster/gizle işlevi eklendi
\par\vspace{3pt}\small $\square$ Ders içinde katılım sergiledim\par $\boxtimes$ Telafi yaptım &
\haftagorsel{hafta3.png}\haftaek{hafta3b.png} \\
\midrule
\textbf{4. hafta} &
Bu haftadaki görevleri sonradan telafi ettim. JavaScript ile DOM manipülasyonu ve olay dinleyicileri (addEventListener) kullanarak etkileşimli bileşenler (artır/azalt sayacı, butona tıklayınca mesaj yazdırma) oluşturdum. Ayrıca HTML formlarından veri alıp (.value) bu verileri ekranda dinamik mesajlara dönüştüren bir uygulama geliştirdim.\newline \textbf{Commit:} JavaScript ile temel DOM manipülasyonu ve olay dinleyici (event listener) örnekleri eklendi
\par\vspace{3pt}\small $\square$ Ders içinde katılım sergiledim\par $\boxtimes$ Telafi yaptım &
\haftagorsel{hafta4.png}\haftaek{hafta4b.png} \\
\midrule
\textbf{5. hafta} &
Bu haftadaki görevleri sonradan telafi ettim. JavaScript dizileri, koşullu ifadeleri (if/else) ve setTimeout fonksiyonunu kullanarak etkileşimli bir "Hafıza Oyunu" (Renk Eşleştirme) kodladım. Kullanıcı tıklamalarını dinleyerek (dataset.renk), eşleşen kartları gizleme (classList.add("gizli")) ve eşleşmeyenleri başlangıç rengine döndürme mantığını kurdum. Ayrıca window.location.href ile buton üzerinden sayfa yönlendirme özelliğini uyguladım.\newline \textbf{Commit:} JavaScript ile etkileşimli Hafıza Oyunu ve sayfa yönlendirme özellikleri eklendi
\par\vspace{3pt}\small $\square$ Ders içinde katılım sergiledim\par $\boxtimes$ Telafi yaptım &
\haftagorsel{hafta5.png}\haftaek{hafta5b.png} \\
\midrule
\textbf{6. hafta} &
Bu haftadaki görevleri sonradan telafi ettim. Harici ve içsel CSS yapılarını kullanarak sayfaların tipografi, renk paleti ve kutu modellemelerini (margin, padding, border) geliştirdim. Ayrıca JavaScript ile onclick olay dinleyicisi ve koşullu yapılar (if/else) kullanarak arayüz üzerinde dinamik renk ve tema geçişleri (gece/gündüz modu) sağlayan etkileşimli bileşenler ekledim.\newline \textbf{Commit:} CSS ile gelişmiş arayüz biçimlendirmesi ve JavaScript tema değiştirme işlevi eklendi
\par\vspace{3pt}\small $\square$ Ders içinde katılım sergiledim\par $\boxtimes$ Telafi yaptım &
\haftagorsel{hafta6.png}\haftaek{hafta6b.png} \\
\midrule
\textbf{7. hafta} &
Bu hafta için bir işleyiş ve etkinlik kaydı bulunmamaktadır.\newline \textbf{Commit:} -
\par\vspace{3pt}\small $\square$ Ders içinde katılım sergiledim\par $\square$ Telafi yaptım &
\haftagorsel{hafta7.png}\haftaek{hafta7b.png} \\
\midrule
\textbf{8. hafta} &
Bu hafta için bir işleyiş ve etkinlik kaydı bulunmamaktadır.\newline \textbf{Commit:} -
\par\vspace{3pt}\small $\square$ Ders içinde katılım sergiledim\par $\square$ Telafi yaptım &
\haftagorsel{hafta8.png}\haftaek{hafta8b.png} \\
\end{longtable}

% ====================================================================
\section{Test ve Değerlendirme}
% ====================================================================
\begin{center}
\rowcolors{2}{acik}{white}
\begin{tabular}{@{}p{6cm}p{3.4cm}p{3.4cm}@{}}
\toprule
\rowcolor{ana}\textcolor{white}{\textbf{Özellik}} & \textcolor{white}{\textbf{Bilgisayar}} & \textcolor{white}{\textbf{Telefon}} \\
\midrule
Anlık Arama ve Filtreleme & Çalışıyor & Çalışıyor \\
Gece Modu ve Metin Boyutu Değişimi & Çalışıyor & Çalışıyor \\
Duyarlı Tasarım (Mobil Kırılımı) & Çalışıyor & Çalışıyor \\
\bottomrule
\end{tabular}
\end{center}

\textbf{Denenen tarayıcılar:} Google Chrome (DevTools simülasyonu ile), Firefox, Safari, Microsoft Edge.

% ====================================================================
\section{Sunum Videosu}
% ====================================================================
\begin{center}
\begin{tabular}{@{}ll@{}}
\textbf{Video süresi:} & 18.00  \\
\textbf{Görünürlük:} & Listelenmemiş \\
\textbf{Bağlantı:} & https://www.youtube.com/watch?v=sknqtlP4rMM \\
\end{tabular}
\end{center}

\subsection{Video Kontrol Listesi}
\begin{itemize}[label={}]
  \item $\boxtimes$ Süre 12--18 dakika aralığında.
  \item $\boxtimes$ Sunum boyunca kameram açık, yüzüm görünüyor.
  \item $\boxtimes$ Görünürlük \textbf{Listelenmemiş} (ya da Herkese açık); \textbf{Özel} değil.
  \item $\boxtimes$ Bağlantıyı gizli pencerede açtım, video oynuyor.
  \item $\boxtimes$ YouTube videoyu işlemeyi bitirdi.
  \item $\boxtimes$ Videoyu baştan sona izledim; ses ve kamera tüm sunumda var.
  \item $\boxtimes$ Kapaktaki bağlantı ile bu bölümdeki bağlantı aynı.
\end{itemize}

% ====================================================================
\section{Sonuç ve Kendi Değerlendirmem}
% ====================================================================
Bu proje, hazır kütüphanelere bağımlı kalmadan saf (Vanilla) HTML, CSS ve JavaScript kullanarak bir ürün geliştirmenin temellerini bana çok net bir şekilde öğretti. Tek bir HTML dosyası kısıtlaması altında çalışmak, başlarda kodların karmaşıklaşacağı endişesi yaratsa da, CSS değişkenlerini doğru kurguladığımda ve HTML semantiğine dikkat ettiğimde "Clean Code" prensiplerinin ne kadar kurtarıcı olduğunu gördüm. En çok zorlandığım nokta, arama ve kategori filtrelemesini aynı fonksiyonda çakışmadan çalıştırmaktı; bu sorunu \texttt{data-} niteliklerini ve mantıksal ve (\&\&) operatörünü birleştirerek çözdüm. İleride bu projeyi geliştirecek olsam, masalların sesli olarak dinlenebilmesi için Web Speech API eklemek ve kullanıcının tema tercihini sayfayı yenilese bile unutmaması için Local Storage sistemini kurmak isterdim.

% ====================================================================
\section{Kaynakça}
% ====================================================================
\begin{thebibliography}{9}
\bibitem{k1} Boratav, Pertev Naili. (2018). \emph{Az Gittik Uz Gittik - Türk Masalları}. Bilgesu Yayıncılık, Ankara. (Nardaniye Hanım ve Sabır Taşı masalları referansı).
\bibitem{k2} Alangu, Tahir. (2001). \emph{Billur Köşk Masalları}. Yapı Kredi Yayınları, İstanbul. (Billur Köşk ve Yedi Başlı Ejderha masalları referansı).
\bibitem{k3} Tezel, Naki. (2001). \emph{Türk Masalları}. Milli Eğitim Bakanlığı Yayınları, İstanbul. (Keloğlan ve Sihirli Tas, Altın Bülbül masalları referansı).
\bibitem{k4} Sakaoğlu, Saim. (2004). \emph{Masal Araştırmaları}. Akçağ Yayınları, Ankara. (Tembel Ahmet ve Üç Turunçlar masalları referansı).
\bibitem{k5} Seyidoğlu, Bilge. (1999). \emph{Erzurum Halk Masalları Üzerine Araştırmalar}. Akçağ Yayınları, Ankara. (Anadolu masalları motif tipolojisi).
\bibitem{k6} MDN Web Docs. \emph{CSS Flexbox ve Vanilla JavaScript Event Listeners}. \url{https://developer.mozilla.org}, erişim tarihi: 3 Kasım 2026.
\bibitem{k7} W3C. \emph{HTML5 Semantic Elements ve Web Erişilebilirliği (A11y)}. \url{https://www.w3.org}, erişim tarihi: 3 Kasım 2026.
\bibitem{k8} Coolors. \emph{Renk Paleti Üretimi ve Kontrast Analizi}. \url{https://coolors.co}, erişim tarihi: 3 Kasım 2026.
\bibitem{k9} Projedeki görseller yapay zeka ile oluşturulmuştur.
\end{thebibliography}

% ====================================================================
\appendix
\section{Teslim Kontrol Listesi}
% ====================================================================
\begin{itemize}[label={}]
  \item $\boxtimes$ MEBİS'e \textbf{yalnızca bu raporun PDF'ini} yükleyeceğim. ZIP ve HTML dosyası yüklemeyeceğim.
  \item $\boxtimes$ PDF dosya adı \texttt{OgrenciNo\_AdSoyad\_vize.pdf} biçiminde.
  \item $\boxtimes$ Kapaktaki GitHub ve YouTube bağlantıları doğru; PDF'te tıklayıp denedim.
  \item $\boxtimes$ GitHub deposunda web sayfası (tek \texttt{.html}), rapor PDF'i, \texttt{.tex} kaynakları ve \texttt{README.md} var.
  \item $\boxtimes$ GitHub deposu \textbf{herkese açık} (ya da öğretim üyesine erişim verdim).
  \item $\boxtimes$ Çalışmamın son hâli teslim saatinden \textbf{önce} depoya gönderildi (sonraki commitler değerlendirilmez).
  \item $\boxtimes$ Sayfa içsel CSS ve içsel JavaScript kullanıyor, harici dosya gerektirmiyor.
  \item $\boxtimes$ Hazır taslak kullanmadım; yapay zekâ kullandıysam Bölüm 5'i doldurdum.
  \item $\boxtimes$ Çalışmanın mantığını sunumda anlatabilirim.
\end{itemize}

\end{document}